\section{从账本看十五世纪的区域负担}

本节按账本的时间序列观察地区、财政、边疆与社会问题。官修实录、诏令和奏疏强于保存中央选择记录的行动与日期；地方志、碑刻、契约、工程遗存及其他政权材料才可能校正具体地点的执行、损失与反应。因此，以下叙事不把行政分类、报送数字或动机判断直接等同于地方社会本身。

\subsection{漕运、灾情与区域行政}

\paragraph{1425（洪熙元年）五月29日：洪熙帝驾崩} 账本首先限定我们能够确认的范围：编年候选的年序可由官修与后出材料交叉；具体月日、参与者、战果或数字必须回查所列材料，不能将单一叙事当作定论。可供互读的材料包括《仁宗实录》洪熙元年条；《明史·本纪》及相关列传；诏令、奏疏或会典版本。这些文本并非同一份现场证言，而是从朝廷、地方或跨政权的位置留下观察。事件之所以进入行政链，与；它带来的可追踪变化是它影响后续人事与政策议程；动机和派系标签不能由单一叙事裁定。译写候选以编年材料定位；实录记载反映朝廷选择，崇祯期须另以奏疏、地方及敌对政权材料交叉。

\paragraph{1425（洪熙元年）六月27日：朱瞻基即位宣德皇帝} 账本首先限定我们能够确认的范围：编年候选的年序可由官修与后出材料交叉；具体月日、参与者、战果或数字必须回查所列材料，不能将单一叙事当作定论。可供互读的材料包括《仁宗实录》洪熙元年条；《明史·本纪》及相关列传；诏令、奏疏或会典版本。这些文本并非同一份现场证言，而是从朝廷、地方或跨政权的位置留下观察。事件之所以进入行政链，与；它带来的可追踪变化是它影响后续人事与政策议程；动机和派系标签不能由单一叙事裁定。译写候选以编年材料定位；实录记载反映朝廷选择，崇祯期须另以奏疏、地方及敌对政权材料交叉。

\paragraph{1425（洪熙元年）九月2日：朱高煦造反} 账本首先限定我们能够确认的范围：编年候选的年序可由官修与后出材料交叉；具体月日、参与者、战果或数字必须回查所列材料，不能将单一叙事当作定论。可供互读的材料包括《仁宗实录》洪熙元年条；《明史·本纪》及相关列传；诏令、奏疏或会典版本。这些文本并非同一份现场证言，而是从朝廷、地方或跨政权的位置留下观察。事件之所以进入行政链，与；它带来的可追踪变化是它影响后续人事与政策议程；动机和派系标签不能由单一叙事裁定。译写候选以编年材料定位；实录记载反映朝廷选择，崇祯期须另以奏疏、地方及敌对政权材料交叉。

\paragraph{1425（洪熙元年）九月22日：朱高煦战败} 账本首先限定我们能够确认的范围：编年候选的年序可由官修与后出材料交叉；具体月日、参与者、战果或数字必须回查所列材料，不能将单一叙事当作定论。可供互读的材料包括《仁宗实录》洪熙元年条；《明史·兵志》及相关列传；边镇、地方志或对方材料。这些文本并非同一份现场证言，而是从朝廷、地方或跨政权的位置留下观察。事件之所以进入行政链，与；它带来的可追踪变化是它改变了后续调兵、补给或地方秩序的条件；战果和伤亡须分开核对。译写候选以编年材料定位；实录记载反映朝廷选择，崇祯期须另以奏疏、地方及敌对政权材料交叉。

\paragraph{1426（宣德元年）：明朝派遣一失哈到野女真建造船坞和仓库} 账本首先限定我们能够确认的范围：编年候选的年序可由官修与后出材料交叉；具体月日、参与者、战果或数字必须回查所列材料，不能将单一叙事当作定论。可供互读的材料包括《宣宗实录》宣德元年条；《明史·兵志》及相关列传；边镇、地方志或对方材料。这些文本并非同一份现场证言，而是从朝廷、地方或跨政权的位置留下观察。事件之所以进入行政链，与；它带来的可追踪变化是它改变了后续调兵、补给或地方秩序的条件；战果和伤亡须分开核对。译写候选以编年材料定位；实录记载反映朝廷选择，崇祯期须另以奏疏、地方及敌对政权材料交叉。

\paragraph{1426（宣德元年）十月5日：林山起义：黎来的军队在二同–秋同战役中对明朝的进攻造成重大伤亡} 账本首先限定我们能够确认的范围：编年候选的年序可由官修与后出材料交叉；具体月日、参与者、战果或数字必须回查所列材料，不能将单一叙事当作定论。可供互读的材料包括《宣宗实录》宣德元年条；《明史·兵志》及相关列传；边镇、地方志或对方材料。这些文本并非同一份现场证言，而是从朝廷、地方或跨政权的位置留下观察。事件之所以进入行政链，与；它带来的可追踪变化是它改变了后续调兵、补给或地方秩序的条件；战果和伤亡须分开核对。译写候选以编年材料定位；实录记载反映朝廷选择，崇祯期须另以奏疏、地方及敌对政权材料交叉。

\subsection{西南、交趾与军费}

\paragraph{1426（宣德元年）冬：林山起义：林山军队将明军驱逐出红河三角洲和越南北部的大部分地区} 账本首先限定我们能够确认的范围：编年候选的年序可由官修与后出材料交叉；具体月日、参与者、战果或数字必须回查所列材料，不能将单一叙事当作定论。可供互读的材料包括《宣宗实录》宣德元年条；《明史·外国传》；《朝鲜王朝实录》、琉球《历代宝案》或海外档案。这些文本并非同一份现场证言，而是从朝廷、地方或跨政权的位置留下观察。事件之所以进入行政链，与；它带来的可追踪变化是它为后续册封、互市或海防安排留下文书与跨政权比较线索。译写候选以编年材料定位；实录记载反映朝廷选择，崇祯期须另以奏疏、地方及敌对政权材料交叉。

\paragraph{1427（宣德二年）十月10日：林山起义：明军援军在林山被围歼} 账本首先限定我们能够确认的范围：编年候选的年序可由官修与后出材料交叉；具体月日、参与者、战果或数字必须回查所列材料，不能将单一叙事当作定论。可供互读的材料包括《宣宗实录》宣德二年条；《明史·兵志》及相关列传；边镇、地方志或对方材料。这些文本并非同一份现场证言，而是从朝廷、地方或跨政权的位置留下观察。事件之所以进入行政链，与；它带来的可追踪变化是它改变了后续调兵、补给或地方秩序的条件；战果和伤亡须分开核对。译写候选以编年材料定位；实录记载反映朝廷选择，崇祯期须另以奏疏、地方及敌对政权材料交叉。

\paragraph{1427（宣德二年）十二月14日：林山起义：明军撤出交趾} 账本首先限定我们能够确认的范围：编年候选的年序可由官修与后出材料交叉；具体月日、参与者、战果或数字必须回查所列材料，不能将单一叙事当作定论。可供互读的材料包括《宣宗实录》宣德二年条；《明史·兵志》及相关列传；边镇、地方志或对方材料。这些文本并非同一份现场证言，而是从朝廷、地方或跨政权的位置留下观察。事件之所以进入行政链，与；它带来的可追踪变化是它改变了后续调兵、补给或地方秩序的条件；战果和伤亡须分开核对。译写候选以编年材料定位；实录记载反映朝廷选择，崇祯期须另以奏疏、地方及敌对政权材料交叉。

\paragraph{1427（宣德二年）：本年朝廷围绕继承、官僚协调和政令传达形成连续记录，可用于核对宣德至正统的中枢运作} 账本首先限定我们能够确认的范围：来源导向的年度桥接条目：本年材料可定位相关政务线索；具体月日、发文机关、地域和执行结果仍须逐项回查，不能以制度文本或奏报替代实际效果。可供互读的材料包括《宣宗实录》宣德二年条；《明史·本纪》及相关列传；诏令、奏疏或会典版本。这些文本并非同一份现场证言，而是从朝廷、地方或跨政权的位置留下观察。事件之所以进入行政链，与；它带来的可追踪变化是它影响后续人事与政策议程；动机和派系标签不能由单一叙事裁定。桥接条目只标示可回查的年度问题与材料链；实录有中央选择性，崇祯期尤其需要奏疏、地方、朝鲜及满文材料交叉。

\paragraph{1428（宣德三年）三月25日：宝藏远航：宣德皇帝命郑和督修大报恩寺} 账本首先限定我们能够确认的范围：编年候选的年序可由官修与后出材料交叉；具体月日、参与者、战果或数字必须回查所列材料，不能将单一叙事当作定论。可供互读的材料包括《宣宗实录》宣德三年条；《明史·本纪》及相关列传；诏令、奏疏或会典版本。这些文本并非同一份现场证言，而是从朝廷、地方或跨政权的位置留下观察。事件之所以进入行政链，与；它带来的可追踪变化是它影响后续人事与政策议程；动机和派系标签不能由单一叙事裁定。译写候选以编年材料定位；实录记载反映朝廷选择，崇祯期须另以奏疏、地方及敌对政权材料交叉。

\paragraph{1428（宣德三年）四月29日：黎黎在后黎王朝的领导下重建了大越王国} 账本首先限定我们能够确认的范围：编年候选的年序可由官修与后出材料交叉；具体月日、参与者、战果或数字必须回查所列材料，不能将单一叙事当作定论。可供互读的材料包括《宣宗实录》宣德三年条；《明史·本纪》及相关列传；诏令、奏疏或会典版本。这些文本并非同一份现场证言，而是从朝廷、地方或跨政权的位置留下观察。事件之所以进入行政链，与；它带来的可追踪变化是它影响后续人事与政策议程；动机和派系标签不能由单一叙事裁定。译写候选以编年材料定位；实录记载反映朝廷选择，崇祯期须另以奏疏、地方及敌对政权材料交叉。

\subsection{北边警报与城市防务}

\paragraph{1428（宣德三年）十月：乌连开袭明边境，宣德皇帝亲自率兵击退} 账本首先限定我们能够确认的范围：编年候选的年序可由官修与后出材料交叉；具体月日、参与者、战果或数字必须回查所列材料，不能将单一叙事当作定论。可供互读的材料包括《宣宗实录》宣德三年条；《明史·兵志》及相关列传；边镇、地方志或对方材料。这些文本并非同一份现场证言，而是从朝廷、地方或跨政权的位置留下观察。事件之所以进入行政链，与；它带来的可追踪变化是它改变了后续调兵、补给或地方秩序的条件；战果和伤亡须分开核对。译写候选以编年材料定位；实录记载反映朝廷选择，崇祯期须另以奏疏、地方及敌对政权材料交叉。

\paragraph{1429（宣德四年）：宣德皇帝在北京郊外进行大检阅} 账本首先限定我们能够确认的范围：编年候选的年序可由官修与后出材料交叉；具体月日、参与者、战果或数字必须回查所列材料，不能将单一叙事当作定论。可供互读的材料包括《宣宗实录》宣德四年条；《明史·兵志》及相关列传；边镇、地方志或对方材料。这些文本并非同一份现场证言，而是从朝廷、地方或跨政权的位置留下观察。事件之所以进入行政链，与；它带来的可追踪变化是它改变了后续调兵、补给或地方秩序的条件；战果和伤亡须分开核对。译写候选以编年材料定位；实录记载反映朝廷选择，崇祯期须另以奏疏、地方及敌对政权材料交叉。

\paragraph{1429（宣德四年）：明军使用携带手炮的骑兵} 账本首先限定我们能够确认的范围：编年候选的年序可由官修与后出材料交叉；具体月日、参与者、战果或数字必须回查所列材料，不能将单一叙事当作定论。可供互读的材料包括《宣宗实录》宣德四年条；《明史·选举志》或《艺文志》；文集、碑刻或书目材料。这些文本并非同一份现场证言，而是从朝廷、地方或跨政权的位置留下观察。事件之所以进入行政链，与；它带来的可追踪变化是它提供制度和传播史的锚点，不能据此推断社会整体接受。译写候选以编年材料定位；实录记载反映朝廷选择，崇祯期须另以奏疏、地方及敌对政权材料交叉。

\paragraph{1429（宣德四年）：漕运、粮饷与军户事务的年度文书显示财政—军事相互约束，定额不能替代实际征解} 账本首先限定我们能够确认的范围：来源导向的年度桥接条目：本年材料可定位相关政务线索；具体月日、发文机关、地域和执行结果仍须逐项回查，不能以制度文本或奏报替代实际效果。可供互读的材料包括《宣宗实录》宣德四年条；《明会典》相关条目；地方志、黄册/契约/河工材料。这些文本并非同一份现场证言，而是从朝廷、地方或跨政权的位置留下观察。事件之所以进入行政链，与；它带来的可追踪变化是其法定安排不等于地方执行，后续须以地域材料追踪负担与成效。桥接条目只标示可回查的年度问题与材料链；实录有中央选择性，崇祯期尤其需要奏疏、地方、朝鲜及满文材料交叉。

\paragraph{1429（宣德四年）：交趾、北边或辽东的军令与边报在本年构成军事决策线索，报称战果应与对方及地方材料比照} 账本首先限定我们能够确认的范围：来源导向的年度桥接条目：本年材料可定位相关政务线索；具体月日、发文机关、地域和执行结果仍须逐项回查，不能以制度文本或奏报替代实际效果。可供互读的材料包括《宣宗实录》宣德四年条；《明史·兵志》及相关列传；边镇、地方志或对方材料。这些文本并非同一份现场证言，而是从朝廷、地方或跨政权的位置留下观察。事件之所以进入行政链，与；它带来的可追踪变化是它改变了后续调兵、补给或地方秩序的条件；战果和伤亡须分开核对。桥接条目只标示可回查的年度问题与材料链；实录有中央选择性，崇祯期尤其需要奏疏、地方、朝鲜及满文材料交叉。

\paragraph{1430（宣德五年）五月：宣德皇帝下令减免所有皇家土地的赋税} 账本首先限定我们能够确认的范围：编年候选的年序可由官修与后出材料交叉；具体月日、参与者、战果或数字必须回查所列材料，不能将单一叙事当作定论。可供互读的材料包括《宣宗实录》宣德五年条；《明会典》相关条目；地方志、黄册/契约/河工材料。这些文本并非同一份现场证言，而是从朝廷、地方或跨政权的位置留下观察。事件之所以进入行政链，与；它带来的可追踪变化是其法定安排不等于地方执行，后续须以地域材料追踪负担与成效。译写候选以编年材料定位；实录记载反映朝廷选择，崇祯期须另以奏疏、地方及敌对政权材料交叉。

\subsection{地方动员与中枢处置}

\paragraph{1430（宣德五年）六月29日：宝藏下西洋：宣德皇帝下令第七次下西洋} 账本首先限定我们能够确认的范围：编年候选的年序可由官修与后出材料交叉；具体月日、参与者、战果或数字必须回查所列材料，不能将单一叙事当作定论。可供互读的材料包括《宣宗实录》宣德五年条；《明史·本纪》及相关列传；诏令、奏疏或会典版本。这些文本并非同一份现场证言，而是从朝廷、地方或跨政权的位置留下观察。事件之所以进入行政链，与；它带来的可追踪变化是它影响后续人事与政策议程；动机和派系标签不能由单一叙事裁定。译写候选以编年材料定位；实录记载反映朝廷选择，崇祯期须另以奏疏、地方及敌对政权材料交叉。

\paragraph{1431（宣德六年）一月19日：宝藏之旅：宝藏船队从南京出发} 账本首先限定我们能够确认的范围：编年候选的年序可由官修与后出材料交叉；具体月日、参与者、战果或数字必须回查所列材料，不能将单一叙事当作定论。可供互读的材料包括《宣宗实录》宣德六年条；《明史·本纪》及相关列传；诏令、奏疏或会典版本。这些文本并非同一份现场证言，而是从朝廷、地方或跨政权的位置留下观察。事件之所以进入行政链，与；它带来的可追踪变化是它影响后续人事与政策议程；动机和派系标签不能由单一叙事裁定。译写候选以编年材料定位；实录记载反映朝廷选择，崇祯期须另以奏疏、地方及敌对政权材料交叉。

\paragraph{1431（宣德六年）三月14日：宝藏之旅：刘家港碑刻立碑} 账本首先限定我们能够确认的范围：编年候选的年序可由官修与后出材料交叉；具体月日、参与者、战果或数字必须回查所列材料，不能将单一叙事当作定论。可供互读的材料包括《宣宗实录》宣德六年条；《明史·本纪》及相关列传；诏令、奏疏或会典版本。这些文本并非同一份现场证言，而是从朝廷、地方或跨政权的位置留下观察。事件之所以进入行政链，与；它带来的可追踪变化是它影响后续人事与政策议程；动机和派系标签不能由单一叙事裁定。译写候选以编年材料定位；实录记载反映朝廷选择，崇祯期须另以奏疏、地方及敌对政权材料交叉。

\paragraph{1431（宣德六年）六月12日：越南黎朝皇帝黎太祖与明朝中国建立朝贡关系，并被明朝皇帝封为安南王} 账本首先限定我们能够确认的范围：编年候选的年序可由官修与后出材料交叉；具体月日、参与者、战果或数字必须回查所列材料，不能将单一叙事当作定论。可供互读的材料包括《宣宗实录》宣德六年条；《明史·外国传》；《朝鲜王朝实录》、琉球《历代宝案》或海外档案。这些文本并非同一份现场证言，而是从朝廷、地方或跨政权的位置留下观察。事件之所以进入行政链，与；它带来的可追踪变化是它为后续册封、互市或海防安排留下文书与跨政权比较线索。译写候选以编年材料定位；实录记载反映朝廷选择，崇祯期须另以奏疏、地方及敌对政权材料交叉。

\paragraph{1431（宣德六年）十二月：宝藏远航：长乐碑刻立，船队从长乐出发} 账本首先限定我们能够确认的范围：编年候选的年序可由官修与后出材料交叉；具体月日、参与者、战果或数字必须回查所列材料，不能将单一叙事当作定论。可供互读的材料包括《宣宗实录》宣德六年条；《明史·本纪》及相关列传；诏令、奏疏或会典版本。这些文本并非同一份现场证言，而是从朝廷、地方或跨政权的位置留下观察。事件之所以进入行政链，与；它带来的可追踪变化是它影响后续人事与政策议程；动机和派系标签不能由单一叙事裁定。译写候选以编年材料定位；实录记载反映朝廷选择，崇祯期须另以奏疏、地方及敌对政权材料交叉。

\paragraph{1432（宣德七年）九月12日：宝藏之旅：宝藏船队抵达萨穆德拉巴赛苏丹国，红宝、马欢随船队访问孟加拉} 账本首先限定我们能够确认的范围：编年候选的年序可由官修与后出材料交叉；具体月日、参与者、战果或数字必须回查所列材料，不能将单一叙事当作定论。可供互读的材料包括《宣宗实录》宣德七年条；《明史·本纪》及相关列传；诏令、奏疏或会典版本。这些文本并非同一份现场证言，而是从朝廷、地方或跨政权的位置留下观察。事件之所以进入行政链，与；它带来的可追踪变化是它影响后续人事与政策议程；动机和派系标签不能由单一叙事裁定。译写候选以编年材料定位；实录记载反映朝廷选择，崇祯期须另以奏疏、地方及敌对政权材料交叉。
